\chapter{Introducere}\label{ch:intro}
\pagestyle{fancy}

Rețelele de calculatoare au devenit infrastructura de bază pe care se sprijină aproape toate activitățile societății moderne, de la comunicații și servicii financiare până la infrastructuri critice și dispozitive interconectate din categoria \textit{Internet of Things} (IoT). Odată cu această dependență crescândă a apărut și o creștere corespunzătoare a numărului și a complexității atacurilor informatice. Atacuri precum refuzul distribuit al serviciului (DDoS), scanările de porturi, atacurile de tip \textit{brute force} sau exfiltrarea de date reprezintă amenințări constante la adresa confidențialității, integrității și disponibilității sistemelor informatice \cite{ref1}. În acest context, mecanismele tradiționale de apărare, cum ar fi firewall-urile și sistemele bazate pe semnături, se dovedesc din ce în ce mai insuficiente, întrucât nu pot detecta atacuri noi, necunoscute anterior, ale căror tipare nu au fost încă înregistrate.

Sistemele de detecție a intruziunilor (\textit{Intrusion Detection Systems}, IDS) reprezintă o linie de apărare esențială în arhitectura de securitate a rețelelor. Un IDS de rețea (\textit{Network Intrusion Detection System}, NIDS) monitorizează traficul care circulă prin rețea și semnalează activitățile suspecte, permițând administratorilor să reacționeze la potențiale incidente de securitate. Spre deosebire de abordările bazate pe semnături, care compară traficul cu o bază de date de tipare cunoscute, abordările bazate pe anomalii încearcă să învețe comportamentul normal al rețelei și să identifice devierile de la acesta. Aceste din urmă abordări au potențialul de a detecta și atacuri necunoscute, motiv pentru care au atras un interes considerabil din partea comunității de cercetare.

În ultimii ani, tehnicile de învățare automată (\textit{machine learning}, ML) s-au impus ca una dintre cele mai promițătoare soluții pentru construirea de sisteme de detecție a intruziunilor bazate pe anomalii. Prin antrenarea unui model de clasificare pe date de trafic etichetate, un NIDS bazat pe ML poate învăța să distingă traficul benign de diferitele tipuri de atac și să generalizeze la situații noi \cite{ref2}. O gamă largă de algoritmi a fost aplicată în acest domeniu, de la clasificatoare liniare simple, precum regresia logistică, până la ansambluri de arbori de decizie, cum sunt \textit{Random Forest} și \textit{Gradient Boosting}, și rețele neuronale profunde. Studiile comparative din literatură arată în mod repetat că modelele de tip ansamblu de arbori obțin rezultate foarte bune pe datele de trafic tabulare, oferind în același timp un grad rezonabil de interpretabilitate \cite{ref3}.

Cu toate acestea, eficacitatea unui NIDS bazat pe ML într-un mediu controlat, de laborator, nu garantează robustețea acestuia în fața unui adversar activ. O problemă din ce în ce mai studiată este cea a atacurilor adversariale: modificări atent construite ale intrării, adesea imperceptibile, care determină un model de învățare automată să producă o clasificare eronată. În contextul securității rețelelor, un atacator care cunoaște, chiar și parțial, comportamentul modelului de detecție poate încerca să modifice ușor traficul malițios astfel încât acesta să fie clasificat drept benign și să treacă nedetectat. Astfel de modificări pot viza caracteristici ale traficului precum valoarea câmpului \textit{Time To Live} (TTL), dimensiunea pachetelor prin \textit{padding}, intervalele dintre pachete sau alte proprietăți care nu afectează funcționalitatea atacului, dar care pot influența decizia modelului. Spre deosebire de domeniul procesării imaginilor, unde atacurile adversariale se bazează în general pe gradientul modelului, în domeniul rețelelor constrângerile suplimentare de validitate a traficului fac ca problema să fie una distinctă și incomplet explorată \cite{ref4}.

Prezenta lucrare se situează la intersecția dintre învățarea automată aplicată securității rețelelor și studiul robusteței modelelor în fața perturbărilor adversariale. Tema abordată constă în evaluarea rezistenței unui model de clasificare a traficului de rețea în fața unor variante ușor modificate ale traficului malițios. Concret, se pornește de la antrenarea unui model de clasificare multi-clasă, capabil să distingă traficul benign de mai multe categorii de atac, pe baza unui set de date public consacrat în domeniul detecției intruziunilor. Ulterior, se generează variante perturbate ale traficului malițios, prin modificări controlate ale unor caracteristici precum TTL sau \textit{padding}, aplicate la niveluri diferite de intensitate. Aceste variante sunt trecute prin model pentru a măsura ce procent dintre ele reușește să evite detecția, precum și pentru a analiza ce tipuri de modificări influențează cel mai mult decizia clasificatorului.

Un aspect particular al abordării multi-clasă îl reprezintă faptul că evaziunea nu se rezumă la situația în care un atac este clasificat drept trafic benign. Un atac perturbat poate rămâne clasificat corect, poate fi clasificat greșit drept trafic normal --- cazul clasic de evaziune, cu potențial ridicat de risc --- sau poate fi confundat cu un alt tip de atac. Această din urmă situație, specifică problemelor multi-clasă, oferă informații suplimentare despre fragilitatea granițelor de decizie ale modelului și este analizată separat în lucrare \cite{ref5}. Distincția între aceste rezultate este importantă din perspectiva securității, întrucât clasificarea eronată a unui atac ca trafic normal are consecințe mult mai grave decât confundarea acestuia cu o altă categorie de atac, care rămâne totuși semnalată ca activitate malițioasă.

Pentru a face rezultatele accesibile și ușor de explorat, lucrarea include și o interfață simplă prin care se pot vizualiza rezultatele și se pot testa diferite niveluri de modificare a datelor. Această componentă permite observarea directă a modului în care intensitatea perturbărilor influențează rata de evaziune și oferă un instrument practic pentru înțelegerea comportamentului modelului în condiții adversariale.

Lucrarea este structurată după cum urmează. Capitolul~\ref{ch:obiective} formulează în detaliu obiectivele proiectului. Capitolul~\ref{ch:stadiul-actual} prezintă stadiul actual al domeniului, incluzând o analiză comparativă a seturilor de date și a modelelor de clasificare disponibile, precum și o trecere în revistă a lucrărilor relevante privind atacurile adversariale asupra sistemelor de detecție a intruziunilor. Capitolul~\ref{ch:fundamentare} descrie fundamentarea teoretică a soluției propuse, iar Capitolul~\ref{ch:proiectare} detaliază proiectarea și implementarea acesteia. Capitolul~\ref{ch:testare} prezintă testarea și validarea, incluzând rezultatele experimentale privind rata de evaziune și influența diferitelor tipuri de perturbări. Capitolul~\ref{ch:manual} conține manualul de instalare și utilizare, iar Capitolul~\ref{ch:concluzii} prezintă concluziile și direcțiile posibile de dezvoltare ulterioară.